% !PRINTED 718
\weberpage{718}
% Source has no running head or folio above "Tabelle IV."; printed 718 is fixed by the
% confirmed sequence 717 (previous page) / 719 (next page).

\begin{center}
\gesperrt{Table} IV.\\
\gesperrt{Developments of the elliptic functions} (p.\ 163).
\end{center}

\weq{1}{\begin{aligned}
\frac{2\varkappa K}{\pi}\operatorname{sn}2Ku &= 4\sin\pi v \sum \frac{(1+q^n)q^{\frac{n}{2}}}{1-2q^n\cos2\pi u+q^{2n}} \\
&= 4\sum \frac{q^{\frac{n}{2}}}{1-q^n}\sin n\pi u.
\end{aligned}} % [sic] printed v in the leading factor; the remaining table uses u

\weq{2}{\begin{aligned}
\frac{2\varkappa K}{\pi}\operatorname{cn}2Ku &= 4\cos\pi v \sum (-1)^{\frac{n-1}{2}} \frac{(1-q^n)q^{\frac{n}{2}}}{1-2q^n\cos2\pi u+q^{2n}} \\
&= 4\sum \frac{q^{\frac{n}{2}}}{1+q^n}\cos n\pi u.
\end{aligned}} % [sic] printed v in the leading factor; the remaining table uses u

\weq{3}{\begin{aligned}
\frac{2K}{\pi}\operatorname{dn}2Ku &= 1+4\sum \frac{(-1)^{\frac{n-1}{2}} q^n(\cos2\pi u-q^n)}{1-2q^n\cos2\pi u+q^{2n}} \\
&= 1+4\sum \frac{q^{\frac{m}{2}}}{1+q^m}\cos m\pi u.
\end{aligned}}

\weq{4}{\begin{aligned}
\frac{2K}{\pi}\frac{1}{\operatorname{sn}2Ku} &= \frac{1}{\sin\pi u} + 4\sin\pi n \sum \frac{q^{\frac{m}{2}}(1+q^m)}{1-q^m\cos2\pi u+q^{2m}} \\
&= \frac{1}{\sin\pi u} + 4\sum \frac{q^n}{1-q^n}\sin\pi n u.
\end{aligned}} % [sic] printed \textquotedbl{}4\sin\pi n\textquotedbl{} and omits the coefficient 2 in the denominator

\weq{5}{\begin{aligned}
\frac{2K}{\pi}\frac{\operatorname{cn}2Ku}{\operatorname{sn}2Ku} &= \operatorname{cotg}\pi u + 4\sin2\pi u \sum \frac{q^{\frac{m}{2}}}{1-q^m\cos2\pi u+q^m} \\
&= \operatorname{cotg}\pi u - 4\sum \frac{q^m}{1+q^m}\sin\pi m u.
\end{aligned}} % [sic] printed numerator q^{m/2} and denominator 1-q^m\cos2\pi u+q^m, expanding the second line gives (-1)^{m/2}q^m over 1-2q^m\cos2\pi u+q^{2m}

\weq{6}{\begin{aligned}
\frac{2K}{\pi}\frac{\operatorname{dn}2Ku}{\operatorname{sn}2Ku} &= \frac{1}{\sin\pi u} + 4\sin\pi u \sum \frac{(-1)^{\frac{m}{2}} q^{\frac{m}{2}}(1-q^m)}{1-q^m\cos2\pi u+q^{2m}} \\
&= \frac{1}{\sin\pi u} - 4\sum \frac{q^n}{1+q^n}\sin n\pi u.
\end{aligned}} % [sic] printed numerator factor (1-q^m) and denominator 1-q^m\cos2\pi u+q^{2m}, expanding the second line gives (1+q^m) and 1-2q^m\cos2\pi u+q^{2m}
